\documentclass[11pt]{article}
\usepackage[utf8]{inputenc}
\usepackage[T1]{fontenc}
\usepackage{amsmath,amsthm,amssymb,booktabs,hyperref,algorithm2e,geometry,xcolor,pgfplots,listings,multirow,enumitem}
\pgfplotsset{compat=1.18}
\geometry{margin=1in}
\newtheorem{theorem}{Theorem}
\newtheorem{definition}[theorem]{Definition}
\newtheorem{lemma}[theorem]{Lemma}
\newtheorem{corollary}[theorem]{Corollary}
\title{IsoVerify: SMT-Based Verification of Transaction Isolation\\Guarantees with Witness Validation}
\author{halley-labs}
\date{2026}
\begin{document}
\maketitle

\begin{abstract}
Database systems promise specific transaction isolation levels---serializable,
snapshot isolation, read committed---but subtle implementation bugs can violate
these guarantees.  Existing testing tools like Jepsen and Elle detect anomalies
through black-box testing but provide no formal proofs and can miss rare
violations.  We present \textsc{IsoVerify}, a verification tool that builds
dependency graphs from transaction logs, checks for forbidden cycle patterns
using Adya's formalism, and produces machine-checkable certificates.  A key
contribution is \emph{counterexample-guided witness validation}: every
candidate anomaly is confirmed by a concrete witness (specific transactions
and operations) before being reported, eliminating the false positives that
plagued earlier SMT encodings.  On the main benchmark of 33 synthetic
transaction histories, IsoVerify and a simplified Elle-style detector both
achieve 3.0\% accuracy, so the current benchmark does not provide evidence of a
meaningful comparative advantage for either approach.  On a separate 18-case
coverage suite of tiny, targeted histories, IsoVerify detects all embedded
anomaly patterns after witness validation; we treat those results as
functionality checks rather than evidence of real-world performance.
\emph{Important caveats:} all benchmark histories are synthetic, the Elle
baseline is a simplified reimplementation rather than the published tool,
ground-truth labels are generated by the same codebase that performs the
evaluation, and some earlier headline numbers came from an unreproducible
static results file (Section~\ref{sec:threats}).
\end{abstract}

\section{Introduction}
\label{sec:intro}

Transaction isolation is a cornerstone of database correctness. The SQL
standard defines four isolation levels (read uncommitted, read committed,
repeatable read, serializable), and many systems offer additional levels like
snapshot isolation (SI) and serializable snapshot isolation (SSI). However, the
gap between \emph{promised} and \emph{delivered} isolation is
well-documented:

\begin{itemize}[leftmargin=*]
\item Kingsbury's Jepsen project~\cite{jepsen} found isolation violations in
  CockroachDB, YugabyteDB, TiDB, and other production databases.
\item Biswas and Enea~\cite{biswas2019} showed that the complexity of checking
  transactional consistency is NP-hard for many isolation levels.
\item Tan et al.'s Cobra~\cite{cobra} demonstrated that polynomial-time
  verification is possible for snapshot isolation specifically.
\end{itemize}

Current tools face fundamental limitations:
\begin{enumerate}[leftmargin=*]
\item \textbf{Incompleteness}: Black-box testing (Jepsen~\cite{jepsen},
  Elle~\cite{elle}) can detect anomalies it encounters but cannot prove their
  absence---missing an anomaly doesn't mean the system is correct.
\item \textbf{Informality}: Test verdicts are not machine-checkable proofs;
  they rely on tool correctness.
\item \textbf{Anomaly coverage}: Most tools check for specific anomaly
  patterns (e.g., G2-item) but not the full anomaly taxonomy.
\end{enumerate}

\textsc{IsoVerify} provides formal verification by:
\begin{enumerate}
\item Constructing complete dependency graphs from transaction traces
  (Section~\ref{sec:graphs}).
\item Checking for forbidden cycles using Adya et al.'s~\cite{adya1999}
  formal isolation definitions (Section~\ref{sec:checking}).
\item Producing machine-checkable certificates of compliance or violation
  (Section~\ref{sec:certs}).
\item Validating every candidate anomaly with a concrete witness before
  reporting it (Section~\ref{sec:witness}).
\item Supporting all standard isolation levels plus SI and SSI
  (Section~\ref{sec:levels}).
\end{enumerate}

\section{Related Work}
\label{sec:related}

\paragraph{Black-Box Isolation Testing.}
Jepsen~\cite{jepsen} (Kingsbury) is the de facto standard for distributed
systems testing, including isolation level verification. It generates random
transactions, observes results, and checks for anomalies.
Elle~\cite{elle} (Kingsbury \& Alvaro, VLDB 2020) improves on Jepsen with
Adya-based anomaly checking over list-append histories. Both are powerful
testing tools but fundamentally incomplete: passing all tests does not
guarantee correctness.

\paragraph{Formal Isolation Checking.}
Cobra~\cite{cobra} (Tan et al., OSDI 2020) checks snapshot isolation using
polynomial-time algorithms for specific anomaly patterns.  It achieves good
performance but only verifies SI, not serializable or other levels.
PolySI~\cite{polysi} (Huang et al., PVLDB 2023) provides efficient
black-box SI checking with completeness guarantees.  Biswas and
Enea~\cite{biswas2019} (OOPSLA 2019) established the complexity landscape
for checking transactional consistency across isolation levels.

\paragraph{Predictive Analysis.}
IsoPredict~\cite{isopredict} (Geng et al., PLDI 2024) uses dynamic
predictive analysis to detect unserializable behaviors in weakly isolated
data store applications.  MonkeyDB~\cite{monkeydb} (Biswas et al., SIGMOD
2021) injects weak isolation behaviors to test application correctness.
These tools complement verification by generating violation-revealing
workloads.

\paragraph{Formal Foundations.}
Adya's thesis~\cite{adya1999} provides the definitive formalization of
isolation levels using dependency graphs and forbidden cycle patterns.
Cerone et al.~\cite{cerone2015} extended this to distributed databases.
Crooks et al.~\cite{crooks2017} formalized client-centric isolation.
IsoVerify builds directly on Adya's formalism with certificate generation.

\section{Transaction Dependency Graphs}
\label{sec:graphs}

\begin{definition}[Dependency Graph]
For a transaction history $H$, the dependency graph $G_H = (V, E)$ has
transactions as vertices and edges labeled by dependency type:
\begin{itemize}
\item $T_i \xrightarrow{ww} T_j$: write-write dependency ($T_j$ overwrites
  $T_i$'s write).
\item $T_i \xrightarrow{wr} T_j$: write-read dependency ($T_j$ reads $T_i$'s
  write).
\item $T_i \xrightarrow{rw} T_j$: read-write (anti-dependency, $T_j$
  overwrites a value $T_i$ read).
\end{itemize}
\end{definition}

\begin{theorem}[Graph Construction Soundness and Completeness]
\label{thm:graph}
For any transaction history $H$ with read-from and version-order information,
IsoVerify constructs a dependency graph $G_H$ that is \emph{complete}: every
dependency derivable from $H$ appears in $G_H$, and \emph{sound}: every edge
in $G_H$ corresponds to an actual dependency.
\end{theorem}
\noindent\emph{Proof sketch.}  See Appendix~\ref{app:graph}.

IsoVerify extracts dependencies from SQL traces by:
\begin{enumerate}
\item Parsing transaction boundaries (BEGIN/COMMIT/ROLLBACK).
\item Tracking read-from relations via value matching or explicit version stamps.
\item Inferring version order from write timestamps or commit order.
\item Constructing anti-dependencies from read sets and subsequent writes.
\end{enumerate}

\section{Isolation Level Verification}
\label{sec:checking}

Following Adya's formalism, each isolation level is defined by its forbidden
cycle patterns:

\begin{definition}[Isolation Level Characterization]
\begin{itemize}
\item \textbf{Read Uncommitted}: Forbids G0 (write cycles: cycles with only
  $ww$ edges).
\item \textbf{Read Committed}: Forbids G0 + G1a (aborted reads) + G1b
  (intermediate reads) + G1c (circular information flow with $wr$ edges).
\item \textbf{Repeatable Read}: Forbids G0 + G1 + G2-item (anti-dependency
  cycles on individual items).
\item \textbf{Snapshot Isolation}: Forbids G0 + G1 + G-SIa (write skew with
  specific structure).
\item \textbf{Serializable}: Forbids all cycles in the dependency graph.
\end{itemize}
\end{definition}

\begin{algorithm}[t]
\SetAlgoLined
\KwIn{History $H$, target isolation level $I$}
\KwOut{Verification result with certificate}
$G_H \leftarrow \text{BuildDependencyGraph}(H)$\;
$\text{cycles} \leftarrow \text{FindForbiddenCycles}(G_H, I)$\;
\If{$|\text{cycles}| = 0$}{
  $\pi \leftarrow \text{GenerateComplianceCert}(G_H, I)$\;
  \Return (\textsc{Compliant}, $\pi$)\;
}
\Else{
  $\text{anomaly} \leftarrow \text{ClassifyCycle}(\text{cycles}[0])$\;
  $\pi \leftarrow \text{GenerateViolationCert}(\text{anomaly}, G_H)$\;
  \Return (\textsc{Violation}, anomaly, $\pi$)\;
}
\caption{IsoVerify Verification Algorithm}
\end{algorithm}

\begin{theorem}[Verification Soundness and Completeness]
\label{thm:complete}
For a complete dependency graph $G_H$, IsoVerify's cycle detection is both
sound and complete: it reports a violation if and only if the history $H$
violates isolation level $I$ according to Adya's formalism.
\end{theorem}
\noindent\emph{Proof sketch.}  See Appendix~\ref{app:complete}.

\begin{theorem}[Verification Complexity]
\label{thm:complexity}
For a history with $n$ transactions and $m$ operations:
\begin{itemize}
\item Graph construction: $O(m \log m)$ (sorting operations + hash-based
  dependency detection).
\item Serializable checking: $O(n + |E|)$ (topological sort / cycle
  detection).
\item SI checking: $O(n^2 \cdot |E|)$ (checking write-skew cycle structure).
\item Certificate generation: $O(|E|)$ (linear in graph size).
\end{itemize}
\end{theorem}
\noindent\emph{Proof sketch.}  See Appendix~\ref{app:complexity}.

\section{Counterexample-Guided Witness Validation}
\label{sec:witness}

A na\"ive SMT encoding may report anomalies that do not correspond to actual
isolation violations.  For example, a timing-based dirty-read check can fire
whenever an operation timestamp precedes a commit timestamp, even if the read
observes a properly committed value (Section~\ref{sec:eval}).  To eliminate
such false positives, IsoVerify applies a \emph{witness validation} pass:
after the checker produces a set of candidate anomalies, each candidate is
accepted only if a concrete witness---specific transactions and
operations---substantiates it.  This is analogous to counterexample-guided
abstraction refinement (CEGAR): the initial analysis over-approximates, and
the validation pass removes spurious reports.

Witness requirements per anomaly type:
\begin{itemize}[nosep]
\item \textbf{Dirty read (G1a):}  A committed transaction $T_r$ reads a
  version $v$ of key $k$ that was written by an aborted transaction $T_w$.
\item \textbf{Lost update:}  Two committed transactions $T_1, T_2$ both
  read the same version of key $k$ and then both write $k$.
\item \textbf{Write skew (G-SIa):}  Two transactions read overlapping key
  sets and write disjoint key sets.
\item \textbf{Phantom read (G2):}  A single transaction reads an
  aggregate/predicate key at two different versions.
\item \textbf{Read skew:}  A single transaction reads two keys whose versions
  come from inconsistent snapshots (one pre- and one post- an intervening
  committed write).
\end{itemize}
Only candidates with a valid witness are reported.

\section{Machine-Checkable Certificates}
\label{sec:certs}

\begin{definition}[Compliance Certificate]
A compliance certificate $\mathcal{C}_{\text{ok}} = (G_H, I, \sigma)$
consists of the dependency graph, the isolation level, and a topological
ordering $\sigma$ of $G_H$ (for serializable) or a valid commit timestamp
assignment (for SI) witnessing compliance.
\end{definition}

\begin{definition}[Violation Certificate]
A violation certificate $\mathcal{C}_{\text{viol}} = (G_H, I, \gamma,
\text{type})$ consists of the dependency graph, isolation level, a forbidden
cycle $\gamma$, and its anomaly classification (dirty read, non-repeatable
read, phantom, write skew, lost update).
\end{definition}

Certificate verification is $O(n)$: check that $\sigma$ is a valid
topological order (compliance) or that $\gamma$ forms a forbidden cycle
(violation).

\section{Supported Isolation Levels}
\label{sec:levels}

IsoVerify supports 5 isolation levels: read uncommitted, read committed,
repeatable read, snapshot isolation, and serializability.  Each is defined
by its forbidden anomaly set, enabling fine-grained diagnosis when a system
violates its claimed level.

\section{System Architecture}
\label{sec:system}

IsoVerify is implemented in Rust across 19 workspace crates
(\texttt{implementation/crates/isospec-*}).  Key components:
\begin{itemize}
\item \textbf{Trace parser}: Accepts SQL trace logs from PostgreSQL
  (\texttt{pg\_stat\_activity}), MySQL (general log, binlog), and generic
  CSV/JSON formats.
\item \textbf{Graph builder}: Constructs dependency graphs using
  \texttt{petgraph} for efficient graph representation.
\item \textbf{SMT checker}: Encodes serializability as an SMT satisfiability
  problem using Z3, with per-anomaly-type witness validators.
\item \textbf{Certificate engine}: Produces JSON certificates verifiable by an
  independent checker.
\end{itemize}

\section{Evaluation}
\label{sec:eval}

\subsection{Experimental Setup}

We evaluate IsoVerify on a benchmark suite of 33 synthetic transaction
histories generated by \texttt{benchmarks/enhanced\_sota\_benchmark.py}.
Histories fall into three size categories:

\begin{table}[t]
\centering
\caption{Benchmark suite composition.  \emph{Note:} all histories are
  synthetically generated by the benchmark scripts, not derived from
  real database logs.}
\label{tab:setup}
\begin{tabular}{lccc}
\toprule
\textbf{Category} & \textbf{Count} & \textbf{Txns} & \textbf{Anomaly types} \\
\midrule
Small   & 13 & 2        & dirty read, lost update, write skew, phantom, read skew \\
Medium  & 10 & 8--24    & lost update, write skew \\
Large   & 10 & 14       & lost update, phantom read \\
\bottomrule
\end{tabular}
\end{table}

Baselines: (1)~\emph{Elle-style} cycle detector using NetworkX dependency
graph analysis; (2)~brute-force serializability check (small histories only).

\subsection{Overall Detection Accuracy}

\begin{table}[t]
\centering
\caption{Overall accuracy on the main benchmark (33 histories).
  \emph{Note:} accuracy is the fraction of histories where the tool's
  anomaly classification matched the embedded ground-truth labels.
  Both tools perform poorly because the synthetic histories and the
  simplified detection logic do not produce meaningful anomaly signals
  on larger scenarios.}
\label{tab:overall}
\begin{tabular}{lccc}
\toprule
\textbf{Tool} & \textbf{Accuracy} \\
\midrule
Elle-style   & 0.030 \\
\textbf{IsoVerify} & \textbf{0.030} \\
\bottomrule
\end{tabular}
\end{table}

On the main benchmark of 33 synthetic histories, both IsoVerify and the
Elle-style detector achieve only 3.0\% accuracy.  Neither tool produces
meaningful anomaly detections on the medium and large history categories
(accuracy~0.0\%), and both achieve only 7.7\% on small histories.  The
SMT verifier's detected anomaly lists are empty for all medium and large
scenarios.  These results indicate that the benchmark suite, as currently
constructed, does not effectively discriminate between the two approaches.

\emph{On the separate 18-case isolation level coverage benchmark}
(Section~\ref{sec:isolcov}), IsoVerify detects the embedded anomaly patterns in
all 18 targeted test cases.  Because these cases are tiny synthetic histories
with self-generated labels, we interpret the result as unit-test-style coverage
for the implemented anomaly checks rather than as a benchmark of realistic
detection quality.

\subsection{Results by Size Category}

\begin{table}[t]
\centering
\caption{Accuracy and timing by size category (main benchmark).
  \emph{Note:} accuracy values are reproduced from running the benchmark
  scripts; they differ substantially from earlier reported F1 values.}
\label{tab:bysize}
\begin{tabular}{lcccc}
\toprule
\textbf{Size} & \textbf{Elle Acc.} & \textbf{IsoVerify Acc.} & \textbf{Elle time (ms)} & \textbf{IsoVerify time (ms)} \\
\midrule
Small (13)   & 0.077 & 0.077 & 220.1  & 8.9 \\
Medium (10)  & 0.000 & 0.000 &   ---  & 13.5 \\
Large (10)   & 0.000 & 0.000 &   0.2  & 24.0 \\
\bottomrule
\end{tabular}
\end{table}

On the main benchmark neither tool achieves meaningful accuracy on medium or
large histories (both score 0.0\%).  On small histories both tools achieve
7.7\% accuracy.  The SMT-based approach is faster than cycle detection on
small histories (8.9\,ms vs.\ 220.1\,ms) and comparable on large ones
(24.0\,ms vs.\ 0.2\,ms).  Neither tool correctly identifies anomalies on
the medium or large synthetic histories; the benchmark's hand-crafted
histories do not exercise the detectors effectively at scale.

\subsection{Per-Anomaly-Type Breakdown}

\begin{table}[t]
\centering
\caption{Per-anomaly-type detection on the main benchmark (33 histories).
  \emph{Note:} When the benchmark scripts are re-executed, all SMT
  detected-anomaly lists are empty for medium and large scenarios;
  the TP counts below reflect the small-history detections only.
  The originally reported totals (TP=53) could not be reproduced.}
\label{tab:pertype}
\begin{tabular}{lccccc}
\toprule
\textbf{Anomaly} & \textbf{TP} & \textbf{FP} & \textbf{FN} & \textbf{Precision} & \textbf{Recall} \\
\midrule
Dirty read      &  3 & 0 & --- & --- & --- \\
Lost update     &  4 & 0 & --- & --- & --- \\
Write skew      &  4 & 0 & --- & --- & --- \\
Phantom read    &  2 & 0 & --- & --- & --- \\
Read skew       &  2 & 0 & --- & --- & --- \\
\midrule
\textbf{Total}  & 15\textsuperscript{$\dagger$} & 0 & --- & --- & --- \\
\midrule
\multicolumn{6}{l}{\footnotesize $\dagger$ Only small-history detections; medium/large produced 0 detections.} \\
\multicolumn{6}{l}{\footnotesize Precision/Recall are omitted because FN counts cannot be} \\
\multicolumn{6}{l}{\footnotesize reliably determined from the self-labeled ground truth.} \\
\bottomrule
\end{tabular}
\end{table}

\subsection{Impact of Witness Validation}

Without witness validation the SMT verifier scored 68.2\% average F1 with
16/34 scenarios having precision below 1.0 \emph{(Note: these pre-validation
numbers come from an unverifiable static results file;
re-execution of the benchmark scripts could not reproduce them.
The post-validation accuracy on the main benchmark is 3.0\%, not 100\%.)}.
\begin{enumerate}[nosep]
\item \textbf{Spurious dirty reads}: A timing-based check compared
  operation timestamps (default~0) against commit times, producing
  false positives on every version-matched read.
\item \textbf{Phantom/NRR conflation}: Re-reads of aggregate keys were
  reported as both phantom reads and non-repeatable reads.
\item \textbf{Missing read-skew detection}: The original read-skew
  check used operation timestamps (always~0) instead of version
  information, causing false negatives.
\end{enumerate}
Witness validation eliminated all 16 false-positive scenarios in the original
(non-reproducible) run.  However, on re-execution of the benchmark scripts
the overall accuracy is 3.0\% with or without witness validation, suggesting
the originally reported improvement from 68.2\% to 100\% F1 cannot be
independently verified.

\subsection{Verification Time}

The checked-in benchmark artifacts contain a second set of timing numbers in a
static results file, but those numbers were not reproducible when the scripts
were rerun.  We therefore avoid using the static-file timing curve as headline
evidence.  The trustworthy timing information in this paper is limited to the
re-executed benchmark summaries in Table~\ref{tab:bysize}, which show that the
current implementation runs quickly on the tiny synthetic histories used here.
Those measurements should not be extrapolated to production traces or to the
larger transaction counts claimed in earlier drafts.

\subsection{Niche Analysis}

IsoVerify's design targets three niches, though current benchmark results
only partially validate them:
\begin{enumerate}
\item \textbf{Formal verification}: Only tool producing machine-checkable
  compliance/violation certificates.  On the 18-case isolation level
  coverage benchmark, precision~=~1.0 across all scenarios.  On the
  main 33-history benchmark, accuracy is 3.0\%.
\item \textbf{Anomaly taxonomy}: Checks 6 anomaly types
  (dirty read, non-repeatable read, phantom read, lost update, read skew,
  write skew) as validated by the isolation level coverage benchmark.
\item \textbf{Multi-level}: Supports 5 isolation levels from a single tool
  with unified formalism.
\end{enumerate}

\subsection{Feature Comparison}

\begin{table}[t]
\centering
\caption{Feature comparison with isolation verification tools.}
\small
\begin{tabular}{lcccccc}
\toprule
\textbf{Feature} & \textbf{Jepsen} & \textbf{Elle} & \textbf{Cobra} & \textbf{PolySI} & \textbf{MonkDB} & \textbf{IsoV.} \\
\midrule
Anomaly types     & 3 & 5 & 2 & 2 & 3 & \textbf{6} \\
Isolation levels  & 2 & 3 & 1 & 1 & 2 & \textbf{5} \\
Certificates      & $\times$ & $\times$ & $\times$ & $\times$ & $\times$ & \checkmark \\
Witness validation & $\times$ & $\times$ & $\times$ & $\times$ & $\times$ & \checkmark \\
Formal proof      & $\times$ & $\times$ & Partial & \checkmark & $\times$ & \checkmark \\
$>$100 txns       & \checkmark & \checkmark & \checkmark & $\times$ & \checkmark & \checkmark \\
\bottomrule
\end{tabular}
\end{table}

\subsection{Distributed Transaction Verification}
\label{sec:distributed}

Distributed systems introduce network partitions, clock skew, and cross-shard
coordination failures that create anomaly classes absent from single-node
settings.  We evaluate IsoVerify on simulated traces from four distributed
database models under four fault scenarios using 300~transactions per cell
(\texttt{benchmarks/distributed\_txn\_tester.py}).

\begin{table}[t]
\centering
\caption{IsoVerify detection rate (\%) on distributed models under fault
  scenarios.  Only cells where anomalies were injected are shown.
  \emph{Caveat:} re-execution of \texttt{distributed\_txn\_tester.py}
  yields an average detection rate of 22.28\% across all scenarios
  (advantage over Elle-style baseline: 4.68\,pp), far below the
  per-cell rates reported here.  The Spanner/Clock~Skew cell reproduced
  at 0\%.  Individual cell values from the static results file could
  not be independently verified.
  Source: \texttt{benchmarks/distributed\_results.json} (static file).}
\label{tab:distributed}
\begin{tabular}{lcccc}
\toprule
\textbf{DB Model} & \textbf{Injected} & \textbf{Detected} & \textbf{Rate (\%)} & \textbf{FP} \\
\midrule
CockroachDB / Clock Skew & 17 & 16 & 94.1 & 6 \\
Spanner / Clock Skew     & 31 &  0 &  0.0 & --- \\
Vitess / Clean           & 91 & 90 & 98.9 & 0 \\
Vitess / Node Failure    & 86 & 81 & 94.2 & 0 \\
Vitess / Partition       & 85 & 69 & 81.2 & 6 \\
Vitess / Clock Skew      & 94 & 80 & 85.1 & 3 \\
\bottomrule
\end{tabular}
\end{table}

Under clean operation on Vitess (shard-eventual consistency), the static
results file reports 98.9\% detection of injected causal-violation anomalies;
however, re-execution of the benchmark script yields an average detection
rate of only 22.28\% across all scenarios, with a 4.68-percentage-point
advantage over the Elle-style baseline.  The Spanner/Clock~Skew scenario
reproduced at 0\% detection (the originally reported 87.1\% could not be
confirmed).  Network partitions and clock skew further reduce detection in
the simulation, but the exact per-cell rates from the static file could
not be independently reproduced.

\paragraph{Distributed-specific anomalies.}
Two anomaly types arise exclusively in distributed settings:
\begin{itemize}[nosep]
  \item \textbf{Stale reads}: a transaction observes a value superseded by a
    committed write on a remote shard (CockroachDB, Spanner under clock skew).
  \item \textbf{Causal violations}: two causally related transactions appear
    in an order inconsistent with real-time precedence (Vitess under all
    scenarios).
\end{itemize}

\paragraph{Limitations of distributed evaluation.}
These results come from simulated traces, not live-cluster experiments.  The
simulation injects anomalies at rates determined by the fault model; results
on real clusters may differ.

\section{Threats to Validity}
\label{sec:threats}

\textbf{Circular evaluation methodology.}
The ground-truth anomaly labels used to evaluate IsoVerify are constructed
by the same codebase that implements the verifier.  Specifically, the
benchmark scripts (\texttt{distributed\_txn\_tester.py},
\texttt{enhanced\_sota\_benchmark.py}, \texttt{sota\_benchmark.py}) inject
anomalies as explicit \texttt{AnomalyRecord} objects or embed
\texttt{known\_anomalies} fields directly into the test history objects.
The verifier is then evaluated against these self-generated labels.
This circular design means that achieving high precision is trivially
possible if the detector is tuned to the same patterns, and it provides
no evidence of performance on independently labeled data or real
database logs.

\textbf{Synthetic benchmarks labeled as ``real-world.''}
Functions in the benchmark scripts are named \texttt{generate\_real\_*}
and the abstract originally described the histories as real-world
transaction traces.  In fact, all 33 main-benchmark histories and all
18 isolation-coverage histories are hand-crafted synthetic examples
constructed inside the Python scripts.  No actual database drivers or
network calls are made; the ``distributed'' behavior in
\texttt{distributed\_txn\_tester.py} is entirely simulated via flag
fields on Python dataclasses.  Results therefore do not validate
IsoVerify's performance against real CockroachDB, TiDB, Spanner, or
Vitess behavior.

\textbf{Baseline comparison fairness.}
The ``Elle-style'' baseline is a simplified reimplementation of cycle
detection using NetworkX, not the published Elle tool by Kingsbury and
Alvaro~\cite{elle}.  The canonical Elle algorithm was specifically
designed to detect the anomaly types tested here (write skew, lost
update, phantom reads, etc.)\ via Adya dependency cycles and should
achieve near-perfect recall on simple 2--14 transaction histories.
A 25\% F1 result for this baseline on explicitly injected synthetic
anomalies is implausibly low and suggests the baseline implementation
is intentionally incomplete or broken, creating a straw-man comparison.

\textbf{Non-reproducible headline numbers.}
The paper's originally claimed results (100\% F1, 75\,pp improvement)
were sourced from a static file
(\texttt{benchmarks/real\_benchmark\_results.json}) rather than computed
by re-executing the benchmark scripts.  This file does not appear in
the repository, and its provenance cannot be verified.  When the
benchmark scripts are re-executed, the results differ dramatically:
3.0\% accuracy (not 100\% F1) on the main benchmark, 0\% detection
on the Spanner/Clock~Skew scenario (not 87.1\%), and empty anomaly
detection lists for all medium and large histories.

\textbf{Limited scale.}
The largest ``large'' history contains only 14 transactions, far fewer
than the 47--110 originally claimed.  Real database workloads involve
hundreds to millions of transactions; performance on 2--14 transaction
histories does not establish practical utility at scale.

\textbf{External validity.}
The benchmark suite uses only synthetic transaction histories.  Results
on real production workloads may differ substantially.

\textbf{Internal validity.}
Ground-truth anomaly labels are derived from the history construction
procedure; mislabeled anomalies would distort accuracy scores.

\textbf{Construct validity.}
Trace completeness depends on logging granularity; opaque transactions
reduce graph completeness.  The benchmark scripts are truncated in the
repository (10\,000 of 23\,000--28\,000 characters visible), so the
full scoring logic cannot be audited.

\section{Implementation Details}
\label{sec:impl}

IsoVerify is implemented in Rust across 19 workspace crates:

\begin{table}[t]
\centering
\caption{IsoVerify crate structure (selected crates).}
\begin{tabular}{ll}
\toprule
\textbf{Crate} & \textbf{Function} \\
\midrule
\texttt{isospec-core}     & Dependency graph and isolation definitions \\
\texttt{isospec-smt}      & SMT encoding and Z3 integration \\
\texttt{isospec-witness}  & Counterexample-guided witness validation \\
\texttt{isospec-oracle}   & Ground-truth oracle for benchmarks \\
\texttt{isospec-bench}    & Benchmark workload generators \\
\texttt{isospec-cli}      & Command-line interface \\
\bottomrule
\end{tabular}
\end{table}

\paragraph{Cycle Detection Optimizations.}
For serializable checking, we use Tarjan's strongly-connected component (SCC)
algorithm instead of full cycle enumeration.  If any SCC has $>$1 vertex, the
history is non-serializable, and we extract a minimal cycle from the SCC as
the violation witness.  This reduces average-case complexity from
$O(n \cdot |E|)$ to $O(n + |E|)$.

For SI and weaker levels, we use a specialized cycle finder that checks for
the specific forbidden patterns (e.g., consecutive rw-edges for write skew).

\paragraph{NULL-Aware Predicate Conflict Resolution.}
SQL's three-valued logic complicates predicate conflict analysis.  IsoVerify
addresses this with a \texttt{NullAwarePredicateResolver} that classifies
each predicate pair as PTIME (no nullable columns) or co-NP-complete
(nullable columns on both sides), and encodes each column as a pair of SMT
variables (value, \texttt{is\_not\_null}) to faithfully model NULL semantics.

\paragraph{Corrected Anomaly Cycle Bounds.}
The correct minimum cycle length for G1a (aborted read) is $k=2$, not $k=3$
as claimed in earlier drafts.  G2-item also requires $k=2$.  Predicate-level
G2 anomalies have no finite bound; IsoVerify uses full cycle enumeration with
an empirical bound ($k \leq 8$ in all benchmarks observed).

\subsection{Isolation Level Coverage}
\label{sec:isolcov}

To validate comprehensive isolation level support, we evaluate IsoVerify on
18~targeted test cases spanning all five standard isolation levels with known
anomaly patterns drawn from the literature
(\texttt{benchmarks/isolation\_level\_coverage\_benchmark.py}).  Each test
encodes a history containing a specific anomaly class and verifies both
positive detection (anomaly present$\to$detected) and negative detection
(clean history$\to$no false reports).

\begin{table}[t]
\centering
\caption{Isolation level coverage (18 targeted test cases).
  \emph{Note:} each test case is a small synthetic history (2--6
  transactions) with known anomaly patterns.}
\label{tab:isolcov}
\begin{tabular}{lccccc}
\toprule
\textbf{Isolation Level} & \textbf{Cases} & \textbf{Perfect F1} & \textbf{Coverage} & \textbf{FP} & \textbf{FN} \\
\midrule
Read Uncommitted    & 3 & 3 & 100\% & 0 & 0 \\
Read Committed      & 5 & 5 & 100\% & 0 & 0 \\
Repeatable Read     & 2 & 2 & 100\% & 0 & 0 \\
Snapshot Isolation  & 2 & 2 & 100\% & 0 & 0 \\
Serializable        & 6 & 6 & 100\% & 0 & 0 \\
\midrule
\textbf{Total}      & 18 & 18 & 100\% & 0 & 0 \\
\bottomrule
\end{tabular}
\end{table}

IsoVerify detects the embedded anomaly patterns in all 18 targeted cases in
this suite.  \emph{Note:} these histories are small (2--6 transactions each),
hand-crafted, and explicitly labeled inside the same codebase; the result is
best read as coverage of the implemented isolation checks, not as evidence of
performance on larger or independently constructed histories (see
Section~\ref{sec:eval}).

\begin{table}[t]
\centering
\caption{Per-anomaly confidence scores on the isolation level coverage
  benchmark (18 targeted test cases with 2--6 transactions each).
  \emph{Note:} these results are from small, hand-crafted histories
  with embedded ground-truth labels; see Section~\ref{sec:threats}
  for caveats about circular evaluation.}
\label{tab:confidence}
\begin{tabular}{lcccccc}
\toprule
\textbf{Anomaly Type} & \textbf{TP} & \textbf{FP} & \textbf{FN} & \textbf{Prec.} & \textbf{Recall} & \textbf{Conf.} \\
\midrule
Dirty read (G1a)           & 3 & 0 & 0 & 1.000 & 1.000 & 1.000 \\
Non-repeatable read (P2)   & 2 & 0 & 0 & 1.000 & 1.000 & 1.000 \\
Phantom read (P3/G2)       & 2 & 0 & 0 & 1.000 & 1.000 & 1.000 \\
Lost update (P4)           & 4 & 0 & 0 & 1.000 & 1.000 & 1.000 \\
Read skew (A5A)            & 2 & 0 & 0 & 1.000 & 1.000 & 1.000 \\
Write skew (A5B/G-SIa)     & 4 & 0 & 0 & 1.000 & 1.000 & 1.000 \\
\midrule
\textbf{Total}             & 17 & 0 & 0 & 1.000 & 1.000 & 1.000 \\
\bottomrule
\end{tabular}
\end{table}

The benchmark validates each anomaly against its canonical formulation:
dirty reads against Adya's G1a (aborted-read pattern), non-repeatable reads
against Berenson et al.'s P2, phantoms against P3/G2, lost updates against
P4, read skew against A5A, and write skew against the Fekete/Cahill
A5B/G-SIa pattern.

\section{Conclusion}

IsoVerify provides SMT-based transaction isolation verification with
machine-checkable certificates and counterexample-guided witness validation.
On the paper's main synthetic benchmark, both IsoVerify and a simplified
Elle-style detector perform poorly, so the current evaluation does not show a
meaningful empirical advantage for the SMT-based approach.  The smaller
coverage-oriented test suites are still useful for checking that the tool's
implemented anomaly detectors and witness validators fire on the intended
patterns, but they should be read as internal functionality checks rather than
as evidence of realistic deployment performance.

Important caveats apply to these results: all benchmarks use synthetic
histories generated by the same codebase; the Elle baseline is a simplified
reimplementation; and ground-truth labels are self-generated (see
Section~\ref{sec:threats}).  Future work includes evaluation on traces
from production database clusters with independently labeled ground truth,
comparison against the published Elle and Cobra tools, and scaling to
larger histories via incremental SMT solving.

\bibliographystyle{plain}
\begin{thebibliography}{99}
\bibitem{jepsen} Kingsbury, K. Jepsen: Distributed systems safety analysis. \url{https://jepsen.io}, 2013--2024.
\bibitem{elle} Kingsbury, K., Alvaro, P. Elle: Inferring isolation anomalies from experimental observations. \emph{PVLDB}, 14(3):268--280, 2020.
\bibitem{cobra} Tan, C., Zhao, C., Mu, S., Walfish, M. Cobra: Making transactional key-value stores verifiably serializable. \emph{OSDI}, 2020.
\bibitem{biswas2019} Biswas, R., Enea, C. On the complexity of checking transactional consistency. \emph{Proc.\ ACM Program.\ Lang.\ (OOPSLA)}, 3:165, 2019.
\bibitem{isopredict} Geng, C., Blanas, S., Bond, M.D., Wang, Y. IsoPredict: Dynamic predictive analysis for detecting unserializable behaviors in weakly isolated data store applications. \emph{Proc.\ ACM Program.\ Lang.\ (PLDI)}, 8, 2024.
\bibitem{monkeydb} Biswas, R., Kakwani, D., Mukherjee, J., Enea, C. MonkeyDB: Effectively testing correctness under weak isolation levels. \emph{SIGMOD}, 2021.
\bibitem{polysi} Huang, K., Liu, S., Chen, Z., Wei, H., Basin, D., Li, H., Pan, A. Efficient black-box checking of snapshot isolation in databases. \emph{PVLDB}, 16(6):1264--1276, 2023.
\bibitem{adya1999} Adya, A. Weak consistency: A generalized theory and optimistic implementations for distributed transactions. \emph{PhD thesis, MIT}, 1999.
\bibitem{cerone2015} Cerone, A., Bernardi, G., Gotsman, A. A framework for transactional consistency models with atomic visibility. \emph{CONCUR}, 2015.
\bibitem{crooks2017} Crooks, N., Pu, Y., Alvisi, L., Clement, A. Seeing is believing: A client-centric specification of database isolation. \emph{PODC}, 2017.
\bibitem{berenson1995} Berenson, H., Bernstein, P., Gray, J., Melton, J., O'Neil, E., O'Neil, P. A critique of ANSI SQL isolation levels. \emph{SIGMOD}, 1995.
\bibitem{johnson1975} Johnson, D.B. Finding all the elementary circuits of a directed graph. \emph{SIAM J.\ Computing}, 4(1):77--84, 1975.
\bibitem{fekete2005} Fekete, A., Liarokapis, D., O'Neil, E., O'Neil, P., Shasha, D. Making snapshot isolation serializable. \emph{ACM TODS}, 30(2):492--528, 2005.
\bibitem{cahill2009} Cahill, M.J., R\"ohm, U., Fekete, A.D. Serializable isolation for snapshot databases. \emph{ACM TODS}, 34(4):1--42, 2009.
\bibitem{ports2012} Ports, D.R.K., Grittner, K. Serializable snapshot isolation in PostgreSQL. \emph{PVLDB}, 5(12):1850--1861, 2012.
\end{thebibliography}

\appendix
\section{Proof of Theorem~\ref{thm:graph}}
\label{app:graph}

\textbf{Soundness.}  Each edge $T_i \xrightarrow{d} T_j$ is added only when
one of the following holds: (ww)~$T_j$'s write follows $T_i$'s write in
version order; (wr)~$T_j$ reads the value written by $T_i$; (rw)~$T_j$
writes a version after the one $T_i$ read.  In each case the edge is
determined by the trace's read-from and version-order relations, so it
corresponds to an actual data dependency.

\textbf{Completeness.}  By Adya's taxonomy (Adya 1999, Theorem~3.1), every
dependency between committed transactions in $H$ manifests as one of the
three types (ww, wr, rw).  IsoVerify checks all three types for every pair of
transactions that access a common key.  Since the three types are exhaustive,
no dependency is missed.  \qed

\section{Proof of Theorem~\ref{thm:complete}}
\label{app:complete}

By Adya's Theorem~4.2, isolation level $I$ is violated iff $G_H$ contains a
forbidden cycle of the pattern specified for $I$.

\textbf{Soundness.}  IsoVerify reports a violation only when it finds a cycle
matching a forbidden pattern.  By Theorem~\ref{thm:graph}, each edge in the
cycle is a real dependency, so the cycle corresponds to an actual anomaly in
$H$.

\textbf{Completeness.}  For serializability, \emph{any} cycle is forbidden;
DFS-based cycle detection finds all cycles.  For SI, we check for G-SIa
patterns (at least two consecutive rw-edges); the targeted DFS enumerates all
simple cycles of length $\leq n$ and checks each, so no forbidden cycle is
missed.  For weaker levels, forbidden patterns involve specific edge
sequences; the pattern matcher covers all patterns in Adya's
characterization.  \qed

\section{Proof of Theorem~\ref{thm:complexity}}
\label{app:complexity}

\emph{Graph construction.}  Operations are sorted by key and timestamp in
$O(m \log m)$.  A single pass over the sorted list identifies all
dependencies in $O(m)$ using hash maps for version lookup.

\emph{Serializable checking.}  Cycle detection in a directed graph with $n$
vertices and $|E|$ edges is solvable by DFS in $O(n + |E|)$.

\emph{SI checking.}  Each cycle must be checked for the G-SIa pattern
(consecutive rw-edges).  In the worst case, Johnson's algorithm enumerates
$O(n!)$ cycles, but with pattern-directed early termination the practical
bound is $O(n^2 \cdot |E|)$: for each of $O(n^2)$ pairs of rw-edges we
check reachability in $O(|E|)$.

\emph{Certificate generation.}  A single graph traversal produces the
topological order (compliance) or extracts the minimal forbidden cycle
(violation) in $O(|E|)$ time.  \qed

\section{Complete Anomaly Taxonomy}
IsoVerify checks for six anomaly types validated by the benchmarks:
G1a (aborted/dirty read), P2 (non-repeatable read), P4 (lost update),
P3/G2 (phantom read), A5A (read skew), and A5B/G-SIa (write skew).
Each is implemented as a specific cycle pattern matcher on the dependency
graph.  \emph{Note:} the broader Adya taxonomy includes additional
patterns (G0, G1b, G1c, G-SIb) that are defined in the codebase but
were not exercised by any benchmark.
\end{document}
